\subsection{The entire nonzero first-coordinate line generated by two}

\begin{statement}{E.52}{Proved}{Complete global sections for every nonzero multiple of the Kummer class of two}
Use the based arithmetic quotient
$H=3\mathbb Z_3\ltimes J$ of E.38 and E.46, where
$J=(3,T-1)\mathbb Z_3[[T-1]]$.
Choose positive compatible real roots of two and put
$\kappa=\kappa_3(2)$.
For a nonzero $r\in\mathbb Z_3$, write uniquely
\[
 r=3^{m-1}e,\qquad m\ge1,\quad e\in\mathbb Z_3^\times,
 \qquad h_m=(3^m-1)/2 .
 \tag{E52.1}
\]
Define
\[
 \begin{split}
 \mathscr P_m^{(2)}
   &=\{p\ne2,3:\ v_3(\operatorname{ord}_{\mathbb F_p^\times}(2))
                                      \le m\},\\
 \mathcal R_m^{(2)}
   &=\{(b_p)_{p\in\mathscr P_m^{(2)}}:
      b_p\in\mathbb Z_3,\ 
      b_p\longrightarrow0\text{ modulo every }3^a\},\\
 \mathcal U_m
   &=\mathbb Z_3^{\,h_m-m}
                       \times\prod_{n\ge m}\mathbb Z_3^{\,3^n}.
 \end{split}
 \tag{E52.2}
\]
The topology on $\mathcal R_m^{(2)}$ is the
inverse-limit topology of its finite-exponent
direct sums. All factors of $\mathcal U_m$ have
their usual compact topology and the displayed
product has the product topology.

The entire set of global section classes of $H$
whose first class is $[3r\kappa]$ has the following
complete coordinates:
\[
 \boxed{\quad
 (b_2,b_3,(b_p),\boldsymbol z_2,\boldsymbol z_3,\boldsymbol u)
 \in
 \mathbb Z_3^2\times\mathcal R_m^{(2)}
                  \times\mathbb Z_3^m\times\mathbb Z_3^m
                  \times\mathcal U_m .
 \quad}
 \tag{E52.3}
\]
The proof specifies norm generators, finite
fundamental-unit matrices, and right inverses
that give an actual homeomorphism. In the scaled
toric target of (E46.7), its value is exactly
\[
 \left(r\kappa,\ 
    b_2\kappa_3(2)+b_3\kappa_3(3)
                   +\sum_{p\in\mathscr P_m^{(2)}}b_p\kappa_3(p)\right).
 \tag{E52.4}
\]
Thus every second toric class supported on
$\{2,3\}\cup\mathscr P_m^{(2)}$ occurs, and its
fiber has the explicit compact coordinates
\[
       \mathbb Z_3^{\,h_m+m}
                       \times\prod_{n\ge m}\mathbb Z_3^{\,3^n}.
 \tag{E52.5}
\]
The finite block in (E52.5) is obtained by the
displayed coordinate concatenation
$(\boldsymbol z_2,\boldsymbol z_3,\boldsymbol u_0)$;
no arithmetic information is discarded.
With origin $g\mapsto(3r\kappa_g,0)$ the entire
fiber (E52.3) is an abelian torsion-free group.
The case $r=0$ is the explicitly
computed fiber (E47.2), related here by the
common projection $H\to3\mathbb Z_3(1)$.

Here are the fields and maps that construct
these coordinates. First take $e=1$, and use the
cofinal levels $n=m+d$, $d\ge0$. Put
\[
 \begin{gathered}
 \alpha_d=2^{1/3^d},\qquad F_d=\mathbb Q(\alpha_d),\\
 \beta_{d,j}=\alpha_d\zeta_{3^{d+j}},\qquad
 L_{d,j}=\mathbb Q(\beta_{d,j})\quad(1\le j\le m),\\
 E_{A,m+d}=F_d\times\prod_{j=1}^m L_{d,j},
       \qquad A=3^m\kappa .
 \end{gathered}
 \tag{E52.6}
\]
The transition maps are cubing of the indicated
generators. Each field has one prime above two
and one above three. At two, their absolute
residue degrees are
\[
               w_0=1,\qquad w_j=\varphi(3^j)\quad(1\le j\le m);
 \tag{E52.7}
\]
at three all residue degrees are one.
Every class number in (E52.6) is prime to three.

In the $j$th complex factor let
$\mathcal E_{d,j}^1$ be the kernel of the
three-primary residue map at two on the
completed ordinary unit group. Then
\[
 N:\mathcal E_{d+1,j}^1\twoheadrightarrow\mathcal E_{d,j}^1,
 \qquad
 \mathcal E_{d,j}^1\simeq\mathbb Z_3^{\,3^{d+j-1}-1}.
 \tag{E52.8}
\]
The real-factor completed unit norm is likewise
onto, with rank $(3^d-1)/2$ at level $d$.
The inverse norm limit of all completed
ordinary units in (E52.6) is precisely
$\mathcal U_m$, by the following explicit
matrix prescription.

Choose fundamental units in each field.
Multiply each complex-factor unit by the
unique power of $\zeta_{3^j}$ which removes
its three-primary residue at two. Their
images are a basis of $\mathcal E_{d,j}^1$.
Together with a fundamental-unit basis in
$F_d$, they give a basis $\boldsymbol\epsilon_d$
of rank
\[
        R_d=(3^{m+d}-1)/2-m .
 \tag{E52.9}
\]
Let $B_d\in M_{R_d,R_{d+1}}(\mathbb Z)$ be
the matrix of the actual relative norms
in these bases, after removing the
prime-to-three torsion. Every Smith invariant
of $B_d$ is a three-adic unit. Choose a
right inverse $S_d$ and a kernel basis $K_d$
over $\mathbb Z_{(3)}$, so that
\[
 B_dS_d=I,\qquad B_dK_d=0,\qquad
 [S_d\ K_d]\in\operatorname{GL}_{R_{d+1}}(\mathbb Z_3).
 \tag{E52.10}
\]
For
$\boldsymbol u=(\boldsymbol u_0,
 \boldsymbol t_m,\boldsymbol t_{m+1},\ldots)\in\mathcal U_m$,
set
\[
 \boldsymbol a_0=\boldsymbol u_0,\qquad
 \boldsymbol a_{d+1}
       =S_d\boldsymbol a_d+K_d\boldsymbol t_{m+d},
 \qquad
 \epsilon(\boldsymbol u)_d
       =\boldsymbol\epsilon_d^{\,\boldsymbol a_d}.
 \tag{E52.11}
\]
This is the promised norm-coherent unit
system, with a unique inverse given by
the same finite matrices.

Choose norm-coherent systems $\xi_{3,j}$
with valuation one along the $j$th
three-prime chain and zero elsewhere,
$0\le j\le m$. The proof constructs them.
The two-prime systems require no choice:
\[
 \xi_{2,0,d}=\alpha_d,\qquad
 \xi_{2,j,d}=\beta_{d,j},
 \tag{E52.12}
\]
placed in the indicated factor and with
one in every other factor. They have
total norms $2^{w_j}$.
One may take, explicitly,
$\xi_{3,0,d}=1+\alpha_d$ and
$\xi_{3,1,d}=1+\beta_{d,1}$, again supported
in their individual factors; their
total norms are three.
For every $p\in\mathscr P_m^{(2)}$, choose
the constructed norm system $\pi_p$
with divisor the unique compatible
degree-one prime chain above $p$.
It has total norm $p$.

Write $\boldsymbol z_i=(z_{i,1},\ldots,z_{i,m})$.
The complete norm-to-section formula for
(E52.3), with $e=1$, is
\[
 \begin{split}
 &\epsilon(\boldsymbol u)\,
 \xi_{2,0}^{\,3b_2-\sum_{j=1}^m w_jz_{2,j}}
       \prod_{j=1}^m\xi_{2,j}^{\,z_{2,j}}\,
 \xi_{3,0}^{\,3b_3-\sum_{j=1}^m z_{3,j}}
       \prod_{j=1}^m\xi_{3,j}^{\,z_{3,j}}\,
 \prod_{p\in\mathscr P_m^{(2)}}\pi_p^{\,3b_p}.
 \end{split}
 \tag{E52.13}
\]
The last product converges in every
finite-exponent quotient. For general $e$,
apply the actual arithmetic automorphism
\[
 \Phi_e:H\longrightarrow H,\qquad
       (a,v(T))\longmapsto(ea,v(T^e)).
 \tag{E52.14}
\]
It preserves augmentation and gives the
first and second coordinates stated in
(E52.4).

The maps from the intrinsic geometric
groups remain the actual quotient maps
\[
 N\twoheadrightarrow H,\qquad
 N/N''\twoheadrightarrow H,\qquad
 N/\overline{\gamma_3(I_Q)}\twoheadrightarrow H
 \tag{E52.15}
\]
with kernels respectively
$[I_y,I_y]$, $[I_y,I_y]/N''$, and
$[I_y,I_y]/\overline{\gamma_3(I_Q)}$
as in (E47.8).
Formula (E52.3) computes the target
section fibers of these maps.
The source lifting obstructions are
retained by those exact kernels.
This addresses the arithmetic section
problem arising from \S51.16,
fols.~733v--736v, scans~172--178,
with the full first-coordinate line
and every continuous norm-compatible
second coordinate specified.
\end{statement}

\begin{proof}
\emph{The actual factorization and its ramification.}
For $A=3^m\kappa$, E.46 gives the finite
etale algebra
$\mathbb Q[X]/(X^{3^{m+d}}-2^{3^m})$.
The exact factorization
\[
 X^{3^{m+d}}-2^{3^m}
  =(X^{3^d}-2)
    \prod_{j=1}^m
       2^{\varphi(3^j)}
          \Phi_{3^j}(X^{3^d}/2)
 \tag{E52.16}
\]
gives (E52.6).
The real factor is Eisenstein at two.
For $j\ge1$, $\beta_{d,j}^{3^d}=2\zeta_{3^j}$
shows that $L_{d,j}$ contains
$K_j=\mathbb Q(\zeta_{3^j})$.
Two is a primitive root modulo $3^j$:
the assertion for $j=1$ and
$v_3(2^{2\cdot3^{j-1}}-1)=j$ prove this
by induction. Thus $K_j$ has one prime
over two, of residue degree
$\varphi(3^j)$.
The valuation of $2\zeta_{3^j}$ there
is one. The equation
$X^{3^d}=2\zeta_{3^j}$ is Eisenstein
at that prime, proving the full degree
$[L_{d,j}:K_j]=3^d$ and total
ramification at two.
This proves irreducibility of every
factor in (E52.16) and all its
specified field embeddings.

The vertical extension
$L_{d+1,j}/L_{d,j}$ is cyclic cubic:
it has degree three by its
two-adic ramification, and the base
contains $\mu_3$.
Its polynomial $X^3-\beta_{d,j}$
shows that only two and three can
ramify. We now prove simultaneously
that three has one prime, totally
ramified in every vertical step,
and that every class number is
prime to three.
At $d=0$ the field is $K_j$,
whose three-class group vanishes
by the one-prime induction proved
in E.47.

Suppose the assertions hold in the
base $B=L_{d,j}$ and put $C=L_{d+1,j}$.
The residue field at two remains
$\mathbb F_{2^{\varphi(3^j)}}$.
In a totally ramified cubic extension
the residue of a unit norm is the
cube of its residue. The reduction
of $\zeta_{3^j}$ has order $3^j$,
and the three-primary order of
this residue multiplicative group
is exactly $3^j$. It is therefore
not a cube. A field element whose
norm is a unit has valuation zero
at this unique prime, so the same
obstruction excludes $\zeta_{3^j}$
from $N_{C/B}C^\times$.

Apply the two exact formulas in
Lemmermeyer's Theorem~1
\cite{E_LemmermeyerAmbiguous}:
\[
 \begin{split}
 \#\operatorname{Am}(C/B)
   &=\frac{h_B3^{t-1}}
       {[\mathcal O_B^\times:
          \mathcal O_B^\times\cap N_{C/B}C^\times]},\\
 \#\operatorname{Am}_{\rm st}(C/B)
   &=\frac{h_B3^{t-1}}
       {[\mathcal O_B^\times:N_{C/B}\mathcal O_C^\times]}.
 \end{split}
 \tag{E52.17}
\]
Here $t\le2$, since the only possible
ramified primes are the unique
primes over two and three, and
there is no real-place ramification.
Every base unit cubed is a norm,
so both displayed indices are
powers of three. The local
obstruction just proved makes
both indices at least three.
Because $3\nmid h_B$, integrality
of the first formula forces $t=2$
and forces its index to equal
three. The three-prime is
therefore ramified and, in this
cubic extension, totally ramified
with one prime above it.
The ambiguous class group has
order prime to three. If the
three-class group of $C$ were
nonzero, its subgroup of
elements killed by three would
be a nonzero vector space in
characteristic three acted on
by a cyclic group of order
three. The nilpotence of
$\sigma-1$ gives a nonzero
fixed vector. That would
give an ambiguous three-class,
which was excluded.
This proves the induction.
The second formula in
(E52.17) also forces its
index to equal three.
Both the element-norm and
the actual unit-norm image
are consequently exactly
the kernel of the nontrivial
residue-cube character at two.

The identical induction applies
to $M_d=K_1F_d$, starting from
$M_0=K_1$, with cubic equation
$X^3-\alpha_d$.
Extension of ideals
$\operatorname{Cl}(F_d)\to
\operatorname{Cl}(M_d)$
followed by norm is
multiplication by two.
It is injective on the
three-primary subgroup,
so $F_d$ has no three-class
group either.
The unique three-prime of
$M_d$, with absolute
ramification index $2\cdot3^d$,
lies over a unique prime of
$F_d$; the intermediate
quadratic extension forces
its ramification index in
$F_d/\mathbb Q$ to be $3^d$.
This proves every assertion
about class numbers and
primes in (E52.6)--(E52.7).

\emph{Actual unit norms, residues, and the free coordinates.}
For a complex factor write
$\mathcal E_{d,j}=
\widehat{\mathcal O_{L_{d,j}}^\times}_3$.
Reduction at two followed by
projection to its three-primary
part gives the split sequence
\[
 0\longrightarrow\mathcal E_{d,j}^1
 \longrightarrow\mathcal E_{d,j}
 \longrightarrow\mu_{3^j}\longrightarrow0.
 \tag{E52.18}
\]
The splitting is the actual
inclusion of $\mu_{3^j}$
into the field.
There are no additional
three-power roots of unity:
reduction at the two-prime
injects such roots into the
fixed residue field, whose
three-primary multiplicative
subgroup has order $3^j$.
It follows that
$\mathcal E_{d,j}^1$ is
torsion-free. Dirichlet's
unit rank and the degree
already proved give its
rank $3^{d+j-1}-1$.

The actual norm index three
in (E52.17) identifies its
image with the residue-cube
subgroup. Every element
of $\mathcal E_{d,j}^1$
therefore has a unit preimage.
The three-primary residue
of any such preimage has
cube one. Multiplication
by the suitable element
of $\mu_3$ removes that
residue and leaves its
norm unchanged, since
$N_{C/B}\zeta_3=\zeta_3^3=1$.
This proves (E52.8).
These deductions hold on
three-completions as stated:
the actual index-three
unit quotient survives
completion exactly, and
all groups of units modulo
torsion are finitely
generated free abelian.

For the real factor,
the three-completed units
of $F_d$ identify with the
complex-conjugation invariants
in those of $M_d$.
Inclusion followed by
the quadratic norm is
multiplication by two,
and the projector is
$(1+c)/2$.
A unit from $F_d$ has
residue one at its
two-prime, whose residue
field is $\mathbb F_2$.
It has a norm preimage
in $M_{d+1}$ by the
already proved unit-norm
criterion. Averaging this
preimage by $(1+c)/2$
gives a preimage in the
completed units of
$F_{d+1}$. Thus these
norm maps are onto
as well.

Under (E52.18), the
map on $\mu_{3^j}$
is cubing, while its
map on the residue-one
units is onto.
An inverse norm system
in the finite residue
group is zero: iterating
$j$ cubings is already
zero. Hence the entire
ordinary-unit norm
limit is exactly the
limit of the real
completed units and
the groups
$\mathcal E_{d,j}^1$.
Their total rank is
\[
 \frac{3^d-1}{2}
      +\sum_{j=1}^m(3^{d+j-1}-1)
       =R_d,\qquad
 R_{d+1}-R_d=3^{m+d}.
 \tag{E52.19}
\]

The normalized fundamental
units in the statement
are explicit bases of
these completed groups.
Their relative norms
are products of the
next base's normalized
fundamental units,
times a root of unity
of order prime to three.
Indeed the remaining
three-power root would
have nontrivial residue.
This defines the
integer matrix $B_d$.
Surjectivity of the
norm proves that all
its nonzero Smith
entries are prime to
three.
For an entirely
specified construction
of (E52.10), take the
lexicographically first
maximal minor whose
determinant is a
three-adic unit.
Its inverse gives
$S_d$ in those pivot
rows, with zero in
the other rows.
For each remaining
column indexed by
$i$, take the kernel
vector $e_i-S_dB_de_i$.
These are the columns
of $K_d$.
The resulting square
matrix $[S_d\ K_d]$
has determinant a
three-adic unit;
its inverse extracts
$B_d\boldsymbol a$
and the nonpivot
coordinates.
This proves (E52.10)
and the bijection
(E52.11), including
its continuous inverse.
All these factors
are compact, so the
inverse-limit
identification has
exactly the product
topology asserted
for $\mathcal U_m$.

\emph{Every surviving valuation, and actual lifting of divisors.}
Outside two and three
the coefficient action
in E.46 is unramified.
For an arithmetic
Frobenius $F_p$, the
action on the position
space $\mathbb Z_3$
is the affine map
\[
           x\longmapsto p x+A_{F_p}.
 \tag{E52.20}
\]
The inertia Tate twist
cancels the coefficient
cyclotomic scalar.
There is an integral
fixed point exactly
when
$x_p=-A_{F_p}/(p-1)$
lies in $\mathbb Z_3$.
For the actual Kummer
torsor of $2^{3^m}$,
this is equivalent
to having compatible
roots of all
three-power orders
in $\mathbb Q_p$.
Cubing is an
automorphism on the
principal units for
$p\ne3$; the finite
residue group
therefore gives
precisely
$p\in\mathscr P_m^{(2)}$.
This argument includes
primes $p\equiv2\pmod3$,
where the fixed point
always exists.

Here is the full
integral measure
calculation behind
this criterion.
When the fixed point
is integral, translate
it to zero and use
the square of the
Frobenius action.
Its action off zero
is multiplication
by $p^2$, and every
nonzero point belongs
to a clopen orbit
homeomorphic to
$\mathbb Z_3$ under
the powers of $p^2$.
A translation-invariant
$\mathbb Z_3$-valued
measure on this
orbit is zero:
the common mass of
a level-$a$ coset
must be divisible
by every power of
three upon refinement.
If no integral
fixed point exists,
then $p\equiv1\pmod3$.
Writing $s=v_3(p-1)$
and $t=v_3(A_{F_p})<s$,
the affine map acts
freely on the
clopen cosets modulo
$3^t$; the identity
\[
 v_3\bigl(\phi^a(x)-x\bigr)
       =t+v_3(a)
 \quad(a\in\mathbb Z_3,\ a\ne0)
 \tag{E52.21}
\]
gives the same orbit
description. Thus
the fixed measure
module is
$\mathbb Z_3\delta_{x_p}$
at an eligible prime
and zero otherwise.
At the finite
positional levels
the orbit sum is
the valuation-one
prime divisor, and
pushforward multiplies
it by the relative
residue degree,
which is the actual
ideal norm.
Consequently the
unique fixed point
is exactly the unique
compatible degree-one
prime chain. All
other valuation
directions have
zero inverse norm
limit.

At two and three
the field calculation
above supplies all
$2(m+1)$ chains,
with relative
residue degree one.
The complete inverse
limit of completed
divisor groups is
therefore
\[
 \mathbb Z_3^{\,m+1}
     \times\mathbb Z_3^{\,m+1}
                    \times\mathcal R_m^{(2)}.
 \tag{E52.22}
\]
The restricted
convergence condition
is necessary already
at positional level
zero and sufficient
at every level:
all fixed-point
measures have
integral unit mass.
This proves the
assertion in the
actual norm system,
including its
inverse-limit
topology.

At every finite
field, its class
group has order
prime to three.
The completed
principal-divisor
map is therefore
onto the completed
divisor group, with
kernel the completed
ordinary units.
Explicitly, a prime
ideal raised to the
class number is
principal, and the
inverse of that
class number in
$\mathbb Z_3$ gives
a completed element
with valuation one.
The kernel assertion
follows by splitting
the original map
onto its free
abelian principal
ideal image and
then taking the
specified
multiplication-by-$3^a$
completions.

For completeness,
the finite residue
unit obstruction is
also removed by
an explicit formula.
Take arbitrary
completed elements
$z_d$ with prescribed
compatible divisors
in one complex
factor, and write
the norm defects
$N z_{d+1}/z_d$ as
units. Let $r_d$
be their
$\mu_{3^j}$ component
in (E52.18). Define
\[
        \rho_d
         =\prod_{a=0}^{j-1}r_{d+a}^{\,3^a}
                \in\mu_{3^j}.
 \tag{E52.23}
\]
Then
$\rho_d=r_d\rho_{d+1}^3$.
Replacing $z_d$
by $z_d\rho_d$
makes every norm
defect have
three-primary
residue one.
The remaining
defects are removed
recursively using
the onto maps
(E52.8), with
the chosen matrix
right inverses.
The real factor
uses its onto
unit norm maps
directly.
This constructs
an actual norm
lift of every
divisor system
in (E52.22).
It also verifies
the inverse-limit
exactness: the
images of ordinary
unit norms stabilize
after at most $m$
steps to their
residue-one subgroups,
and these subgroups
have onto norm
transitions.
No uncomputed
$\varprojlim^1$
term is omitted.

The specified
$\xi_{3,j}$ and
$\pi_p$ now follow
by applying this
construction to
their individual
prime chains.
For the other
generators use
the identities
\[
 \begin{gathered}
 N_{d+1/d}(\alpha_{d+1})=\alpha_d,\qquad
 N_{d+1/d}(\beta_{d+1,j})=\beta_{d,j},\\
 N_{d+1/d}(1+\alpha_{d+1})=1+\alpha_d,\qquad
 N_{d+1/d}(1+\beta_{d+1,1})=1+\beta_{d,1}.
 \end{gathered}
 \tag{E52.24}
\]
Each equality is
the exact norm in
an odd cubic
extension, computed
from $X^3-a$.
The total norms
of $1+\alpha_d$
and $1+\beta_{d,1}$
are respectively
$3$ and
$N_{K_1/\mathbb Q}(1+2\zeta_3)=3$.
Their only
nonzero valuation
is therefore one
at the indicated
three-prime.
This proves all
the literal
generator claims.

\emph{The full section coordinates and their topology.}
Divide an arbitrary
norm system by
the systems just
constructed with
its actual prime
valuations. The
quotient is an
ordinary-unit
system and has
the unique form
$\epsilon(\boldsymbol u)$
by (E52.11).
Thus the valuation
and unit coordinates
are both exhaustive
and unique.
If $d_{2,j}$ and
$d_{3,j}$ denote
the bad-prime
valuations, its
level-zero norm
has valuation
$\sum_jw_jd_{2,j}$
at two,
$\sum_jd_{3,j}$
at three, and
$c_p$ at an
eligible prime.
The exact cube
condition of E.46
is therefore
\[
 \sum_{j=0}^m w_jd_{2,j}\in3\mathbb Z_3,\qquad
 \sum_{j=0}^m d_{3,j}\in3\mathbb Z_3,\qquad
 c_p\in3\mathbb Z_3.
 \tag{E52.25}
\]
Solving these
three displayed
conditions with
the coefficient
one of $d_{2,0}$
and $d_{3,0}$
gives exactly
(E52.13).
The infinite
product converges
because at every
coefficient exponent
only finitely
many $b_p$ are
nonzero. Valuations
give a continuous
inverse to these
coordinates; the
remaining unit
inverse is the
compact continuous
matrix inverse
(E52.11).
This proves the
homeomorphism,
its toric formula,
and all its
stated fibers.
Changes of the
chosen prime
norm lifts
multiply them
by definite
ordinary-unit
norm systems,
and change
$\boldsymbol u$
by the corresponding
convergent linear
series. They
preserve the
displayed toric
map.

Finally,
substitution
$T\mapsto T^e$
is a continuous
automorphism of
$\Lambda$, with
inverse
$T\mapsto T^{e^{-1}}$.
It preserves $J$
and intertwines
$T^a$ with
$T^{ea}$.
This verifies
the group law
in (E52.14).
The arithmetic
action is
$g(a,v)=(\chi_g a,
\chi_g v(T^{\chi_g}))$;
these substitutions
commute with
$\Phi_e$.
Augmentation
is unchanged.
Hence
$\Phi_e$ gives
the exact
transport from
the computed
first cocycle
$3^m\kappa$
to $3r\kappa$,
and changes
no second
toric coordinate.
All groups in
(E52.3) are
torsion-free,
which proves
the final group
assertion.

\emph{Executed finite certificates.}
The retained programs
\texttt{e\_two\_units.gp} and
\texttt{e\_two\_line.py}
record the actual integral bases,
fundamental units, normalized units,
relative norms, and matrix splittings.
The PARI computation uses
\texttt{bnfinit(P,1)} followed by
\texttt{bnfcertify}; the latter
returns one for each of the eight
fields
\[
 F_0,F_1,F_2,\quad
 L_{0,1},L_{1,1},L_{2,1},\quad
 L_{0,2},L_{1,2}.
 \tag{E52.26}
\]
Thus every finite unit basis and
class group used in these checks
is certified without a hypothesis
on zeros of zeta functions
\cite{E_PARICertifiedUnits}.
All eight certified class numbers
are one. In the recorded normalized
bases the three norm matrices
with nonzero target rank are
\[
 \begin{gathered}
 B_{F_2/F_1}=\begin{pmatrix}3&2&-1&1\end{pmatrix},\\
 B_{L_{2,1}/L_{1,1}}=
 \begin{pmatrix}
 0&-3&0&0&0&0&1&1\\
 3&-3&1&1&0&0&1&1
 \end{pmatrix},\\
 B_{L_{1,2}/L_{0,2}}=
 \begin{pmatrix}
 0&-3&0&1&1&0&1&-1\\
 3&-3&1&0&0&1&0&-1
 \end{pmatrix}.
 \end{gathered}
 \tag{E52.27}
\]
Their Smith entries are respectively
$(1)$, $(1,1)$, and $(1,1)$.
The exact rational right inverses,
kernel bases, and inverse coordinate
matrices are retained in the JSON
certificate. The Python program
independently computes all twenty-three
relative unit norms as multiplication
determinants in the actual cubic
algebras, verifies their unit
relations and all matrix identities,
and checks the rank formulas,
prime-chain criteria, residue correction,
and complete valuation coordinates.
These finite certificates verify the
specified maps; the preceding
induction and norm-limit arguments
prove the assertions at every level.
\end{proof}
